% Transcrição/tradução independente completa a partir das pp. físicas 6--8 do IAS,
% pp. impressas 1--3. Nenhum componente NomiL ou inglês mantido contém este texto.

\chapter*{Um fio de Ariadne para o SGA~4, o SGA~$4\frac{1}{2}$ e o SGA~5}
\addcontentsline{toc}{chapter}{Um fio de Ariadne para o SGA 4, o SGA 4 1/2 e o SGA 5}

As exposições I a VI do SGA~4 apresentam a teoria geral das topologias de
Grothendieck. Por serem muito detalhadas, podem revelar-se valiosas no estudo
de topologias exóticas, como a que dá origem à cohomologia cristalina. Para a
topologia étale, tão próxima da intuição clássica, não é necessário um aparato
de segurança tão imponente: basta conhecer, por exemplo, o livro de
Godement~[4] e ter um pouco de fé. Outras referências possíveis são os
capítulos I a III das notas de Artin~[1], a exposição Bourbaki de Giraud~[3]
ou [Arcata]~I. As exposições VII e VIII iniciam o estudo da topologia étale.
São mais detalhadas que o capítulo III de Artin e que [Arcata]~II.

O estudo das curvas é a chave da cohomologia étale. Artin o inicia na
Exposição~IX do SGA~4; na Exposição~X são apresentadas consequências no que
diz respeito à dimensão cohomológica. Os pontos essenciais são retomados em
[Arcata]~III. A Exposição~XI não é necessária para o que segue; contém uma
elegante demonstração do teorema de comparação entre a cohomologia étale e a
cohomologia clássica no caso particular das variedades algébricas complexas
lisas.

% SGA-SRC-000064 (aceito, apenas evidência em português): o impresso traz ``la
% dimension cohomologiques'' na p. física 6. O português usa o singular idiomático
% ``dimensão cohomológica'' nos dois modos e não fabrica um erro.
As exposições XII e XIII demonstram o teorema fundamental de mudança de base
para um morfismo próprio. Como observou Artin em
\emph{Algebraic approximation of structures over complete local rings},
Publ. Math. IHÉS 36 (1969), pp.~23--58, \S~3, seu teorema de aproximação
permite simplificar a demonstração. Esse caminho é seguido em [Arcata]~IV,
onde os múltiplos \emph{dévissages} da demonstração mal são esboçados e as
aplicações da Exposição~XIV são indicadas apenas muito brevemente. O teorema
permite definir a ``cohomologia com suportes compactos'' e os funtores
``imagens diretas superiores com suportes próprios'' $R^i f_!$ (cujas
propriedades se reduzem às dos funtores de imagens diretas superiores para um
morfismo próprio). Esse é o objeto da prolixa Exposição~XVII; o essencial é
dito em [Arcata]~IV, embora algumas partes da Exposição~XVII possam ter
utilidade independente.

As exposições XV e XVI se concentram no teorema fundamental de aciclicidade
local para os morfismos lisos. A influência do SGA~7 faz-se sentir na
exposição paralela, menos detalhada, [Arcata]~V, que contém também a
demonstração dos resultados do SGA~2 que são necessários.

A exposição [Arcata]~VI é dedicada à dualidade de Poincaré; é mais clara que o
SGA~4, Exposição~XVIII, mas não pretende apresentar uma demonstração completa.
Para quem deseja uma demonstração puramente algébrica, esta é simplificada por
uma referência a [Dualité]~\S~3; ela não contém a demonstração do formalismo ao
qual obedece o morfismo de traço (XVIII, \S~2).

Na cohomologia étale existe um formalismo de dualidade análogo ao da
cohomologia coerente. Para estabelecê-lo, Grothendieck utilizava a resolução
de singularidades e a conjectura de pureza (quanto ao enunciado, veja-se
[Cycle]~2.1.4), estabelecida em um contexto relativo no SGA~4, Exposição~XVI,
e ---módulo a resolução--- em característica igual no SGA~4,
Exposição~XIX. Os pontos-chave são estabelecidos por outro método em
[Th. finitude], para os esquemas de tipo finito sobre um esquema regular de
dimensão 0 ou 1. No SGA~5, Exposição~I, são apresentados diversos
desenvolvimentos. No SGA~5, Exposição~III, mostra-se como esse formalismo
implica a fórmula dos traços de Lefschetz-Verdier em sua forma muito geral.

Vê-se que, na versão original do SGA~5, a fórmula de Lefschetz-Verdier era
estabelecida apenas conjecturalmente. Além disso, ali não eram calculados os
termos locais. Para a aplicação às funções $L$, este seminário contém outra
demonstração ---esta, sim, completa--- no caso particular do morfismo de
Frobenius: a que se encontra em [Rapport]. Outras referências: para o
enunciado e o esquema dos \emph{dévissages}, a exposição Bourbaki de
Grothendieck~[5]; para uma breve descrição da redução, devida a Grothendieck,
do caso crucial a um caso já tratado por Weil, [2], \S~10; para um tratamento
$\ell$-ádico deste último caso, [Cycle], \S~3.

No SGA~5 encontra-se também um tratamento detalhado da passagem da cohomologia
com coeficientes em $\mathbf{Z}/\ell^n$ à cohomologia com coeficientes em
$\mathbf{Z}_\ell$ (exposições V e VI, por J.~P.~Jouanolou), bem como do
formalismo do morfismo de Frobenius (Exposição~XV, por C.~Houzel). As
exposições X e XII, redigidas por I.~Bucur, apresentam o cálculo da
característica de Euler-Poincaré de um feixe sobre uma curva e o cálculo do
traço do endomorfismo de sua cohomologia definido por certas correspondências.
Na Exposição~VII, Jouanolou apresenta o cálculo da cohomologia de esquemas
clássicos e algumas aplicações. Por outro lado, a exposição de Serre
``Introduction à la théorie de Brauer'' foi retomada em seu livro~[6]
(terceira parte).

\section*{Bibliografia}

\begin{description}
\item[{[1]}] M.~Artin, \emph{Grothendieck topologies}, Harvard University,
1962.

\item[{[2]}] P.~Deligne, \emph{Les constantes des équations fonctionnelles
des fonctions L}, em \emph{Modular functions of one variable II},
pp.~501--597, Lecture Notes in Math.~349, Springer-Verlag.

\item[{[3]}] J.~Giraud, \emph{Analysis situs}, Séminaire Bourbaki 256,
maio de 1963, Benjamin.

\item[{[4]}] R.~Godement, \emph{Topologie algébrique et théorie des
faisceaux}, Act. Sci. et Ind. 1252, Hermann, 1958.

\item[{[5]}] A.~Grothendieck, \emph{Formule de Lefschetz et rationalité des
fonctions L}, Séminaire Bourbaki 279, dez. 1964, Benjamin.

\item[{[6]}] J.~P.~Serre, \emph{Représentations linéaires des groupes finis},
Collection Méthodes, Hermann, 1967.
\end{description}

\begin{description}
\item[SGA] \emph{Séminaire de Géométrie algébrique du Bois-Marie}.

\item[SGA 4] \emph{Théorie des topos et cohomologie étale des schémas}
(dirigido por M.~Artin, A.~Grothendieck e J.~L.~Verdier), Lecture Notes 269,
270 e 305.

\item[SGA 5] \emph{Cohomologie $\ell$-adique et fonctions L} (dirigido por
A.~Grothendieck), de publicação próxima em Lecture Notes.
\end{description}
